\begin{abstract}
\textbf{Background} Accurate classification and segmentation of brain tumors in MRI and CT scans are essential for diagnosis and treatment planning. However, the heterogeneous morphology of brain tumors, including irregular shapes, sizes, and spatial variability, makes this task highly challenging. Traditional convolutional neural networks (CNNs) lack rotational and translational invariance, which limits their ability to generalize across different orientations.

\textbf{Methods} This study introduces a geometric deep learning framework called Modified Special Euclidean (Mod-SE(2)), which integrates geometric priors to enhance spatial consistency and reduce reliance on data augmentation. We extend this framework with a \textbf{2.5D Spatial Context Model} that stacks adjacent slices ($n-1, n, n+1$) to capture vital 3D spatial depth while retaining the computational efficiency of 2D networks. By incorporating symmetry-preserving group convolutions and spatial priors via an \textbf{SE(2)-Equivariant CNN (SE2-CNNET)}, our architecture structurally guarantees rotational equivariance. To further delineate ambiguous tumor boundaries, we introduce an advanced \textbf{Edge-Boundary Loss} function utilizing morphological operations to heavily penalize boundary misclassifications. Mod-SE(2) was evaluated on multiple medical image datasets for classification (Mod-Cls-SE(2)) and segmentation (Mod-Seg-SE(2)) tasks.

\textbf{Results} Mod-Cls-SE(2) achieved an average classification accuracy of 0.914, outperforming ResNet101, VGG16, and their variants. For segmentation tasks, extensive 10-fold cross-validation demonstrates that our approach significantly outperforms standard U-Net baselines in both spatial awareness and boundary precision. The model achieved a high Dice coefficient and an IoU exceeding standard benchmarks, notably reducing false positive predictions at tumor margins.

\textbf{Conclusions} The upgraded Mod-SE(2) leverages 2.5D geometric priors to improve spatial consistency, efficiency, and interpretability in brain tumor analysis. Its symmetry-aware design enables better generalization across tumor shapes and outperforms traditional methods across key metrics. The Edge-Boundary Loss ensures accurate boundary delineation, supporting neurosurgical planning and downstream applications.
\end{abstract}
